\documentclass[11pt]{article}

% This standalone wrapper is intended for internal review.  The four input
% files can also be copied directly into an ACS manuscript or Supporting
% Information document.  Build and retain only the final PDF with:
%   latexmk -pdf manuscript_simulation_theory.tex
%   latexmk -c manuscript_simulation_theory.tex
\usepackage[margin=1in]{geometry}
\usepackage[T1]{fontenc}
\usepackage{lmodern}
\usepackage{microtype}
\usepackage{amsmath,amssymb,mathtools,bm}
\usepackage[version=4]{mhchem}
\usepackage{siunitx}
\usepackage{booktabs}
\usepackage{enumitem}
\usepackage[numbers,sort&compress]{natbib}
\usepackage[colorlinks=true,allcolors=blue]{hyperref}

\newcommand{\kOn}{k_{\mathrm{on}}}
\newcommand{\kOff}{k_{\mathrm{off}}}
\newcommand{\tauB}{\tau_{B,\mathrm{eff}}}
\newcommand{\NA}{N_{\mathrm A}}
\newcommand{\dt}{\Delta t}
\newcommand{\dx}{a}
\newcommand{\Prb}{\mathop{\mathrm{Pr}}}
\newcommand{\E}{\mathop{\mathrm{E}}}

\title{Mathematical Description of the Single-Ligand\\
Escape and Rebinding Simulations}
\author{Methods text supporting manuscript Figures 2--5}
\date{\today}

\begin{document}
\maketitle

\begin{abstract}
This document defines the physical and mathematical model used for the
single-ligand escape and rebinding simulations in manuscript Figures 2--5.
The model is an on-lattice, particle-based reaction--diffusion Monte Carlo
scheme.  Diffusion is derived from Fick's equation and the Einstein relation,
whereas receptor association and dissociation are represented by Poisson
reaction hazards.  Single-ligand trajectories measure the cumulative time
spent bound before escape from an explicit near-surface region.  A
renewal--reward balance then converts that effective bound time into an
occupancy-aware concentration response and an operational $C_{50}$.  This
kinetic $C_{50}$ is deliberately distinguished from the thermodynamic
dissociation constant $K_D$.
\end{abstract}

\section{Scope and model classification}

The simulation resolves one tagged ligand moving on a cubic lattice above a
voxelized sensor surface.  It is most precisely described as an
\emph{on-lattice reaction--diffusion Monte Carlo} or single-particle,
reaction--diffusion-master-equation-like model.  It is Brownian in the
overdamped, long-time sense because its propagator has the Fickian mean-square
displacement, but it is not an inertial molecular-dynamics or underdamped
Langevin simulation.  This class of spatial stochastic model is closely
related to lattice reaction--diffusion methods and particle-based stochastic
reaction--diffusion algorithms used for biochemical systems
\citep{gillespie1976,andrews2004,isaacson2008,hellander2012}.

The model separates two questions.  First, a tagged analyte that begins bound
to a receptor is followed through dissociation, free diffusion, and any
subsequent self- or cross-rebinding events until it escapes or reaches the
observation horizon.  Second, the resulting cumulative bound time is inserted
into a mean-field kinetic balance to predict the concentration-dependent
fractional occupancy.  No claim is made that nanostructuring changes the
microscopic equilibrium free energy or the intrinsic $K_D=\kOff/\kOn$.

\begin{table}[ht]
\centering
\caption{Principal symbols.  Concentration is expressed in molar units
throughout the kinetic equations.}
\label{tab:notation}
\begin{tabular}{@{}lll@{}}
\toprule
Symbol & Meaning & Units \\
\midrule
$a$ & cubic lattice spacing & m \\
$\dt$ & diffusion/reaction update interval & s \\
$D$ & ligand diffusion coefficient & \si{\meter\squared\per\second} \\
$V_{\mathrm{rxn}}$ & local reaction volume & L \\
$\rho_R$ & receptor surface density & \si{\per\meter\squared} \\
$\kOn$ & intrinsic bimolecular association constant & \si{\liter\per\mole\per\second} \\
$\kOff$ & unimolecular dissociation constant & \si{\per\second} \\
$\tauB(\phi)$ & mean cumulative bound time before escape & s \\
$\phi$ & imposed background receptor occupancy & dimensionless \\
$\theta$ & predicted ensemble fractional occupancy & dimensionless \\
$C$ & bulk analyte concentration & M \\
\bottomrule
\end{tabular}
\end{table}

\begin{table}[ht]
\centering
\caption{Shared manuscript simulation settings unless explicitly swept.  The
geometry, receptor density, $\kOn$, $\kOff$, and background occupancy vary in
the corresponding factorial and kinetic studies.}
\label{tab:baseline-settings}
\begin{tabular}{@{}lll@{}}
\toprule
Setting & Baseline value & Implementation \\
\midrule
Nominal domain & $100\times100\times100$ nm$^3$ & explicit near-surface region \\
Lattice spacing & $a=\SI{1}{\nano\meter}$ & cubic lattice \\
Diffusion coefficient & $D=\SI{1e-11}{\meter\squared\per\second}$ & isotropic \\
Reaction region & $n_{\mathrm{rxn}}=1$ voxel & $V_{\mathrm{rxn}}=a^3$ \\
Association constant & $\kOn=10^7$ M$^{-1}$ s$^{-1}$ & baseline chemistry \\
Dissociation constant & $\kOff=10^{-2}$ s$^{-1}$ & baseline chemistry \\
Lateral boundaries & periodic in $x$ and $y$ & tiled sensor patch \\
Escape boundary & $z_{\mathrm{esc}}=\SI{75}{\nano\meter}$ & absorbing plane \\
Initial state & receptor-bound & one tagged ligand \\
Trajectories & 100 per release receptor & per independent replicate \\
Baseline horizon & $t_{\max}=10^4$ s & right censoring at horizon \\
Kinetic-sweep horizon & $t_{\max}=100/\kOff$ & constant intrinsic-lifetime multiple \\
\bottomrule
\end{tabular}
\end{table}

\section{Spatial discretization and physical time}
\label{sec:grid}

\subsection{Cubic lattice and voxelized sensor surface}

The explicit domain has nominal dimensions $L_x\times L_y\times H$ and is
partitioned into a simple cubic lattice with spacing $a$.  The implemented
numbers of lattice intervals are
\begin{equation}
 N_x=\operatorname{round}\!\left(\frac{L_x}{a}\right),\qquad
 N_y=\operatorname{round}\!\left(\frac{L_y}{a}\right),\qquad
 N_z=\operatorname{round}\!\left(\frac{H}{a}\right).
 \label{eq:grid-counts}
\end{equation}
Fluid and solid sites are stored on the resulting Cartesian grid.  An
analytical surface---planar, cylindrically curved, or sinusoidal---is converted
to a solid mask.  Every face separating an accessible fluid voxel from a
solid voxel is an exposed surface face with area $a^2$.  Consequently, the
voxelized reactive area is
\begin{equation}
 A_s=\sum_{f\in\mathcal F_{\mathrm{reactive}}}a^2,
 \label{eq:voxel-area}
\end{equation}
where $\mathcal F_{\mathrm{reactive}}$ contains the bulk-accessible exposed
faces.  At prescribed receptor density $\rho_R$, the receptor count is
\begin{equation}
 N_R=\operatorname{round}(\rho_R A_s),
 \label{eq:receptor-count}
\end{equation}
unless a fixed receptor count is explicitly imposed.  Receptor-bearing faces
are sampled uniformly without replacement for the manuscript simulations;
there is therefore at most one receptor per surface face.  Fixed-$N_R$
controls separate geometric confinement from changes in total receptor
number.

The staircase representation in Eq.~\ref{eq:voxel-area} is part of the
numerical model.  Refining $a$ improves geometric resolution but also changes
the local reaction volume and, in a local lattice reaction model, can change
the relationship between microscopic and mesoscopic reaction rates.  Spatial
resolution must therefore be assessed by convergence studies rather than by
assuming that arbitrarily small voxels are automatically more physical
\citep{isaacson2008,hellander2012,hellander2015}.

\subsection{Diffusive update interval}

The coefficient $\lambda_D$ is not an additional physical parameter; it is
the dimensionless group produced by discretizing Fick's second law.  Let
$c_{i,j,k}^{n}$ denote the concentration at lattice site $(i,j,k)$ and time
$t_n=n\dt$.  A forward difference in time and centered second differences in
space give
\begin{align}
 \frac{c_{i,j,k}^{n+1}-c_{i,j,k}^{n}}{\dt}
 =\frac{D}{a^2}\big(&c_{i+1,j,k}^{n}+c_{i-1,j,k}^{n}
 +c_{i,j+1,k}^{n}+c_{i,j-1,k}^{n}\nonumber\\
 &+c_{i,j,k+1}^{n}+c_{i,j,k-1}^{n}
 -6c_{i,j,k}^{n}\big).
 \label{eq:ftcs-grid-derivation}
\end{align}
Multiplying Eq.~\ref{eq:ftcs-grid-derivation} by $\dt$ shows that the fraction
transferred from a site to \emph{each} nearest neighbor during one update is
$D\dt/a^2$.  This is the standard forward-time, centered-space (FTCS)
discretization of the diffusion equation \citep{crank1975}.  We therefore
define
\begin{equation}
 \lambda_D=\frac{D\dt}{a^2}.
 \label{eq:lambda-D}
\end{equation}
It can also be read as the ratio between the update duration $\dt$ and the
diffusive timescale $a^2/D$ associated with one lattice spacing.  In the
single-particle representation, the ligand moves to any of its six nearest
neighbors with probability $\lambda_D$ per direction and remains at its
present site with probability
\begin{equation}
 p_0=1-6\lambda_D.
 \label{eq:stay-probability}
\end{equation}
Nonnegative probabilities require the standard three-dimensional explicit
diffusion stability condition
\begin{equation}
 \dt\leq \frac{a^2}{6D}.
 \label{eq:stability}
\end{equation}
When the user does not prescribe a timestep, the implementation uses
\begin{equation}
 \dt=0.95\frac{a^2}{6D},
 \label{eq:default-dt}
\end{equation}
which gives $p_0=0.05$ and $p_{\pm x}=p_{\pm y}=p_{\pm z}=0.95/6$.  For the
manuscript baseline $a=\SI{1}{\nano\meter}$ and
$D=\SI{1e-11}{\meter\squared\per\second}$,
Eq.~\ref{eq:default-dt} gives $\dt\approx\SI{1.58e-8}{\second}$.  Physical
observation times are obtained as integer numbers of these updates.

\subsection{Numerical scales that must remain distinct}

Three lengths have different meanings and should not be conflated:
\begin{enumerate}[leftmargin=*]
 \item $a$ controls the spatial resolution, surface voxelization, and local
 reaction volume;
 \item the geometric length scales (for example, a cylindrical radius or
 sinusoidal wavelength) define confinement and curvature; and
 \item the escape height defines the boundary of the explicitly simulated
 near-surface region and hence the first-passage endpoint.
\end{enumerate}
The supplementary numerical controls vary $a$, lateral domain size, and
escape height because each can otherwise introduce an arbitrary computational
scale into a rebinding calculation.

\section{Brownian diffusion and the escape observable}
\label{sec:diffusion}

\subsection{From Fick's law to a lattice random walk}

In a homogeneous fluid without advection or external forces, conservation of
mass together with Fick's first law, $\bm J=-D\bm\nabla c$, gives Fick's second
law
\begin{equation}
 \frac{\partial c}{\partial t}=D\nabla^2c.
 \label{eq:fick2}
\end{equation}
The forward-Euler, centered-space discretization of Eq.~\ref{eq:fick2} at
lattice site $\bm i$ is
\begin{equation}
 c_{\bm i}^{n+1}
 =c_{\bm i}^{n}
 +\lambda_D\sum_{\bm e\in\mathcal N_6}
 \left(c_{\bm i+\bm e}^{n}-c_{\bm i}^{n}\right),
 \label{eq:ftcs}
\end{equation}
where $\mathcal N_6=\{\pm\hat{\bm x},\pm\hat{\bm y},\pm\hat{\bm z}\}$.
Equation~\ref{eq:ftcs} is also the master equation of the random-walk kernel
defined by Eqs.~\ref{eq:lambda-D} and \ref{eq:stay-probability}.  Thus, the
single-particle Monte Carlo propagator and the finite-difference diffusion
equation have the same ensemble evolution.

For one update, symmetry gives $\E[\Delta\bm r]=\bm 0$ and
\begin{equation}
 \E\!\left[\lvert\Delta\bm r\rvert^2\right]
 =6\lambda_Da^2=6D\dt.
 \label{eq:one-step-msd}
\end{equation}
Independent increments therefore give
\begin{equation}
 \E\!\left[\lvert\bm r(t)-\bm r(0)\rvert^2\right]=6Dt,
 \qquad
 \E[(x(t)-x(0))^2]=2Dt,
 \label{eq:einstein-msd}
\end{equation}
which is the Einstein relation for three-dimensional Brownian diffusion
\citep{einstein1905}.  The lattice walk consequently represents overdamped
Brownian transport on times and lengths resolved by $\dt$ and $a$.

\subsection{Boundary conditions}

The manuscript simulations use periodic lateral boundaries,
\begin{equation}
 (x,y,z)\equiv(x+L_x,y,z)\equiv(x,y+L_y,z),
 \label{eq:periodic-boundaries}
\end{equation}
implemented by modular wrapping of the $x$ and $y$ indices.  This suppresses
artificial loss from the sides and represents a periodically tiled sensor
patch.  A move proposed into the voxelized solid is rejected, which is the
random-walk realization of a reflecting, no-flux boundary on nonreactive
surface encounters.  Receptors on exposed faces add stochastic reaction to
this otherwise reflecting interface.

For the escape assay, a trajectory terminates when the free ligand reaches an
absorbing escape surface.  The principal manuscript calculations use a
horizontal plane $z=z_{\mathrm{esc}}$; crossing an explicitly open domain
boundary also counts as escape.  If
\begin{equation}
 T_{\mathrm{esc}}=\inf\{t>0:z(t)\ge z_{\mathrm{esc}}\},
 \label{eq:first-passage}
\end{equation}
then the retention survival function is the first-passage quantity
\begin{equation}
 S(t)=\Prb(T_{\mathrm{esc}}>t).
 \label{eq:survival}
\end{equation}
It measures persistence in the explicit near-surface region, not the survival
of a chemical bond.  A trajectory can alternate many times between bound and
unbound states before the event in Eq.~\ref{eq:first-passage}.

Trajectories that have not escaped by the finite horizon $t_{\max}$ are
right-censored.  With $d_j$ escape events among $n_j$ trajectories at risk at
ordered event time $t_j$, the plotted Kaplan--Meier estimator is
\begin{equation}
 \widehat S(t)=\prod_{t_j\le t}\left(1-\frac{d_j}{n_j}\right).
 \label{eq:kaplan-meier}
\end{equation}
The assay additionally records the observation time, number of binding
intervals, self- and cross-rebinding counts, and cumulative bound time
\begin{equation}
 T_B=\sum_{q=0}^{N_{\mathrm{reb}}}B_q,
 \label{eq:cumulative-bound-time}
\end{equation}
where $B_0$ is the initial bound interval and subsequent $B_q$ are rebound
intervals.  Free-excursion durations contribute to $T_{\mathrm{esc}}$ but not
to $T_B$.  This distinction is essential because receptor occupancy is a
time-integrated bound-state observable.

\section{Association, dissociation, and rebinding}
\label{sec:binding}

In stochastic kinetics, a \emph{hazard} is the instantaneous conditional rate
of an event, given that the event has not already occurred.  If $T$ is the
waiting time to an event, its hazard is
\begin{equation}
 h(t)=\lim_{\delta t\rightarrow0}
 \frac{\Prb\!\left(t\leq T<t+\delta t\mid T\geq t\right)}{\delta t}.
 \label{eq:hazard-definition}
\end{equation}
It has units of inverse time.  For a constant hazard $h$, the probability of
no event during an interval of length $t$ is $e^{-ht}$, and the probability of
at least one event is $1-e^{-ht}$.  Thus, ``hazard'' here simply means the
instantaneous stochastic event rate; it does not imply a separate physical
mechanism.  This rate-to-probability conversion is the same one used in
chemical master-equation and kinetic Monte Carlo methods
\citep{gillespie1976}.

\subsection{Microscopic association probability}

The chemical step is
\begin{equation}
 \ce{L + R <=>[k_{on}][k_{off}] LR},
 \label{eq:reaction}
\end{equation}
where $\kOn$ is supplied in \si{\liter\per\mole\per\second} (equivalently,
$\mathrm{M}^{-1}\,\mathrm{s}^{-1}$) and $\kOff$ in
\si{\per\second}.  A ligand can react when it occupies the fluid voxel adjacent
to a receptor-bearing surface face.  If the configured reaction region
contains $n_{\mathrm{rxn}}$ voxels, its volume in liters is
\begin{equation}
 V_{\mathrm{rxn}}
 =n_{\mathrm{rxn}}a^3\left(10^3\ \mathrm{L\,m^{-3}}\right).
 \label{eq:reaction-volume}
\end{equation}
One ligand in this volume corresponds to molar concentration
\begin{equation}
 C_1=\frac{1}{\NA V_{\mathrm{rxn}}},
 \label{eq:one-particle-concentration}
\end{equation}
where $\NA$ is Avogadro's constant.  For $m$ available receptors at the same
reaction site, elementary mass action gives the total instantaneous
association rate (hazard)
\begin{equation}
 h_{\mathrm{on}}(m)=m\kOn C_1
 =\frac{m\kOn}{\NA V_{\mathrm{rxn}}}.
 \label{eq:on-hazard}
\end{equation}
Treating reaction events as a Poisson process during one update yields
\begin{equation}
 p_{\mathrm{bind}}(m)
 =1-\exp[-h_{\mathrm{on}}(m)\dt]
 =1-\exp\!\left[-\frac{m\kOn\dt}{\NA V_{\mathrm{rxn}}}\right].
 \label{eq:binding-probability}
\end{equation}
Conditional on association, one of the $m$ available receptors is selected
uniformly.  The exponential form preserves $0\le p_{\mathrm{bind}}\le1$ and
reduces to the familiar chemical-master-equation propensity
$p_{\mathrm{bind}}\simeq m\kOn\dt/(\NA V_{\mathrm{rxn}})$ when the exponent is
small \citep{gillespie1976,fange2006}.

Equation~\ref{eq:binding-probability} defines how the reported intrinsic
$\kOn$ is interpreted at the chosen lattice scale.  In a diffusion-influenced
reaction, the macroscopic association rate also depends on transport and on
the chosen reactive length scale; this distinction is familiar from
Smoluchowski and radiation-boundary descriptions
\citep{collins1949,agmon1990,andrews2009}.  Moreover, a strictly local
reaction--diffusion master equation does not possess a mesh-independent
microscopic limit for every bimolecular system
\citep{isaacson2008,hellander2012}.  Accordingly, $a$ and
$n_{\mathrm{rxn}}$ are declared model parameters, and lattice-spacing
robustness is tested explicitly rather than absorbed into an unreported rate
renormalization.

\subsection{Dissociation probability and event-wise acceleration}

For a bound complex with first-order dissociation constant $\kOff$, the
continuous waiting-time density is
\begin{equation}
 f_B(t)=\kOff e^{-\kOff t},
 \qquad \E[B]=\frac{1}{\kOff}.
 \label{eq:exponential-dwell}
\end{equation}
The dissociation hazard is the constant first-order rate $h_{\mathrm{off}}=
\kOff$.  Integrating this rate over one update gives the exact conditional
dissociation probability
\begin{equation}
 p_{\mathrm{off}}=1-e^{-\kOff\dt}.
 \label{eq:off-probability}
\end{equation}
Rather than evaluating Eq.~\ref{eq:off-probability} repeatedly during long
bound intervals, the escape assay draws the number of bound updates directly,
\begin{equation}
 G\sim\operatorname{Geometric}(p_{\mathrm{off}}),\qquad B_{\Delta}=G\dt,
 \label{eq:geometric-dwell}
\end{equation}
with support $G=1,2,\ldots$.  This event-wise skip is statistically identical
to testing dissociation after every lattice update.  Its discrete mean is
\begin{equation}
 \E[B_{\Delta}]
 =\frac{\dt}{1-e^{-\kOff\dt}}
 =\frac{1}{\kOff}+\frac{\dt}{2}+O(\kOff\dt^2),
 \label{eq:discrete-dwell-mean}
\end{equation}
so the physical exponential mean is recovered as $\kOff\dt\rightarrow0$.
This acceleration changes computational cost but not the discrete stochastic
model.

Each receptor-conditioned escape trajectory begins in the bound state.  Its
first interval is therefore drawn from Eq.~\ref{eq:geometric-dwell} and
contributes to $T_B$, but it is not counted as a rebinding event.  Upon
dissociation, the ligand is placed in the fluid voxel adjacent to the released
receptor.  It subsequently diffuses according to Section~\ref{sec:diffusion}
and can bind again according to Eq.~\ref{eq:binding-probability}.

\subsection{Self-rebinding, cross-rebinding, and finite occupancy}

A new binding event is classified relative to the receptor vacated at the
start of the current free excursion.  Return to that receptor is
\emph{self-rebinding}; binding to a different receptor is
\emph{cross-rebinding}.  The all-rebinding model permits both.  A mechanistic
self-only control disables cross-rebinding while leaving diffusion,
$\kOn$, and $\kOff$ unchanged.

Finite ensemble occupancy is represented without adding other explicitly
diffusing ligands.  For a prescribed background occupancy $\phi$, every
non-source receptor is independently available with probability $1-\phi$
during a tagged ligand's free excursion.  That availability field remains
fixed throughout the excursion and is resampled after the next dissociation;
it therefore does not flicker on the diffusion-update timescale.  The receptor
just vacated by the tagged ligand is forced available, because its ligand is
the tagged molecule itself.  This construction preserves self-return while
progressively suppressing cross-rebinding as the background approaches
saturation.

The background field is a controlled mean-field approximation, not an
explicit many-ligand equilibrium simulation.  It assumes statistically
independent receptor occupancy outside the tagged trajectory and neglects
ligand--ligand interactions and correlated turnover of background ligands.
Its advantage is that it isolates how receptor availability modifies the
single-ligand retention kernel at a known occupancy.

\section{From cumulative bound time to fractional occupancy}
\label{sec:occupancy}

\subsection{Classical two-state reference}

For independent receptors exposed to a spatially uniform bulk concentration
$C$, classical mass-action kinetics gives
\begin{equation}
 \frac{d\theta}{dt}=\kOn C(1-\theta)-\kOff\theta.
 \label{eq:langmuir-kinetics}
\end{equation}
At steady state,
\begin{equation}
 \theta_{\mathrm{classical}}(C)
 =\frac{\kOn C}{\kOn C+\kOff}
 =\frac{C}{C+K_D},
 \qquad K_D=\frac{\kOff}{\kOn}.
 \label{eq:langmuir-isotherm}
\end{equation}
This is the noncooperative Langmuir reference for a one-to-one binding site
\citep{langmuir1918}.  It assumes that dissociation ends the occupied episode;
the same ligand's subsequent return is not retained as memory in the
two-state rate equation.

\subsection{Renewal--reward closure with rebinding}

Let
\begin{equation}
 \tauB(\phi)=\E[T_B\mid\text{background occupancy }\phi]
 \label{eq:effective-bound-time}
\end{equation}
be the mean cumulative time for which a tagged ligand is bound before it
escapes the near-surface region.  It includes the initial bound interval and
all subsequent bound intervals, but excludes intervening free-diffusion time.
Thus, $\tauB$ is the total bound-state ``reward'' generated by one successful
arrival into the retained trajectory.

For $N_R$ statistically equivalent receptor sites with ensemble occupancy
$\theta$, let $\lambda_{\mathrm{in}}$ denote the total rate at which ligands
associate with vacant receptors and initiate retained trajectories.  Under
the mean-field association model,
\begin{equation}
 \lambda_{\mathrm{in}}
 =N_R\kOn C(1-\theta).
 \label{eq:arrival-rate}
\end{equation}
The expected number of bound receptors at steady state is the trajectory
initiation rate multiplied by the mean cumulative bound time generated by
each trajectory,
\begin{equation*}
 \left\langle N_B\right\rangle
 =\lambda_{\mathrm{in}}\tauB(\theta).
\end{equation*}
By the definition of fractional occupancy, the same expected number is
\begin{equation*}
 \left\langle N_B\right\rangle=N_R\theta.
\end{equation*}
Combining these two identities with Eq.~\ref{eq:arrival-rate} gives
\begin{equation*}
 N_R\theta=N_R\kOn C(1-\theta)\tauB(\theta).
\end{equation*}
Cancelling $N_R$ gives the occupancy balance
\begin{equation}
 \theta=\kOn C(1-\theta)\tauB(\theta).
 \label{eq:renewal-balance}
\end{equation}
For an occupancy-independent effective lifetime, Eq.~\ref{eq:renewal-balance}
can be written
\begin{equation}
 \theta(C)=\frac{\kOn C\tauB}{1+\kOn C\tauB}.
 \label{eq:effective-occupancy}
\end{equation}
This is equivalently a two-state kinetic expression with operational removal
rate $k_{\mathrm{off,eff}}=1/\tauB$.  In the continuous-time no-rebinding
limit, $\tauB=1/\kOff$, and Eq.~\ref{eq:effective-occupancy} reduces to the
classical result in Eq.~\ref{eq:langmuir-isotherm}.  In the implemented
discrete process, the corresponding no-rebinding mean is given by
Eq.~\ref{eq:discrete-dwell-mean} and differs from $1/\kOff$ by a relative
error of order $\kOff\dt$.  This error is negligible for the manuscript
timestep and is independently testable by timestep refinement.

At finite occupancy, $\tauB$ is not constant because occupied neighboring
receptors block cross-rebinding.  The simulations therefore impose
$\phi\in(0,1)$, estimate $\tauB(\phi)$, and identify the self-consistent
ensemble occupancy as $\theta=\phi$.  Rearrangement of
Eq.~\ref{eq:renewal-balance} gives the parametric concentration-response curve
used in Figures~3--5:
\begin{equation}
 \boxed{
 C(\phi)=
 \frac{\phi}{(1-\phi)\kOn\tauB(\phi)}
 }
 \qquad (0<\phi<1).
 \label{eq:concentration-from-phi}
\end{equation}
The curve is therefore not an extrapolation of a single dilute-limit lifetime;
each point uses a retention simulation performed at the corresponding
background occupancy.

\subsection{Operational half-occupancy concentration and normalization}

The concentration producing a chosen occupancy $p$ is
\begin{equation}
 C_p=\frac{p}{(1-p)\kOn\tauB(p)}.
 \label{eq:Cp}
\end{equation}
In particular,
\begin{equation}
 \boxed{
 C_{50}=\frac{1}{\kOn\tauB(0.5)}
 },
 \qquad
 \boxed{
 \frac{C_{50}}{K_D}=\frac{1}{\kOff\tauB(0.5)}
 }.
 \label{eq:C50}
\end{equation}
The dimensionless residence amplification
\begin{equation}
 A_B(\phi)=\kOff\tauB(\phi)
 \label{eq:residence-amplification}
\end{equation}
equals unity in the classical, no-rebinding limit.  More generally,
normalizing Eq.~\ref{eq:concentration-from-phi} by
$K_D=\kOff/\kOn$ gives
\begin{equation}
 \frac{C(\phi)}{K_D}
 =\frac{\phi}{(1-\phi)\kOff\tauB(\phi)}
 =\frac{\phi}{(1-\phi)A_B(\phi)}.
 \label{eq:normalized-concentration}
\end{equation}
At the same occupancy, the concentration required by the classical model is
$C_{p,\mathrm{classical}}=K_Dp/(1-p)$.  Hence the kinetic concentration
reduction is
\begin{equation}
 \frac{C_{p,\mathrm{classical}}}{C_p}
 =\kOff\tauB(p)=A_B(p).
 \label{eq:concentration-reduction}
\end{equation}
This identity makes clear why occupancy-dependent suppression of
cross-rebinding can change the response-curve shape: if $A_B(p)$ is larger at
low $p$ than near saturation, the dilute end is shifted left more strongly
than the saturated end.

\subsection{Estimating concentration from occupancy}

If the measured fractional occupancy $\theta$ and the appropriate
occupancy-conditioned lifetime are known, the model's concentration estimate
is the inverse of Eq.~\ref{eq:effective-occupancy}:
\begin{equation}
 \boxed{
 \widehat C
 =\frac{\theta}{(1-\theta)\kOn\tauB(\theta)}
 }.
 \label{eq:estimated-concentration}
\end{equation}
This estimator is model-based: it requires the same intrinsic $\kOn$,
geometry, receptor density, and occupancy convention used to calibrate
$\tauB(\theta)$.  Near $\theta=0$ or $1$, uncertainty is amplified by the
logit factor $\theta/(1-\theta)$, so endpoint estimates require correspondingly
larger trajectory ensembles.

For receptor-resolved maps, the same equations are applied to
$\tau_{B,r}(\phi)$, the mean cumulative bound time of trajectories initiated
at receptor $r$.  This produces local operational quantities such as
\begin{equation}
 \frac{C_{50,r}}{K_D}
 =\frac{1}{\kOff\tau_{B,r}(0.5)},
 \label{eq:local-C50}
\end{equation}
which report spatially heterogeneous retention on a single structured
surface.  They are local kinetic response metrics, not receptor-specific
binding free energies.

\subsection{Censoring and statistical interpretation}

For each receptor location and condition, $\tauB$ is estimated by averaging
$T_B$ over independent trajectories and then combining independent simulation
replicates.  If a trajectory reaches $t_{\max}$ before escape, its recorded
$T_B$ is a lower bound on the eventual cumulative bound time.  The simple
sample mean used in the concentration conversion is therefore a restricted
mean and is biased downward if censoring is appreciable; this would bias
$C_{50}$ upward.  Manuscript production runs set the horizon far beyond the
intrinsic mean bound time (and kinetic sweeps scale it with $1/\kOff$), while
the censoring fraction is retained as a quality-control observable.  The
Kaplan--Meier estimator in Eq.~\ref{eq:kaplan-meier} is used for the separate
first-passage survival plots.

Equations~\ref{eq:renewal-balance}--\ref{eq:estimated-concentration} also assume
that bulk concentration is spatially equilibrated and that arrival to vacant
receptors follows the pseudo-first-order rate $\kOn C$.  The single-ligand
escape assay resolves post-dissociation retention; it does not independently
calculate macroscopic mass-transfer limitations from a finite reservoir.


\section{Connection to manuscript Figures 2--5}

The mathematical objects above enter the figures as follows.  Figure~2 uses
the first-passage survival function, the distribution of cumulative bound
time, and the number of rebindings to establish how curvature and receptor
density alter retention.  Figure~3 uses occupancy-aware values of
$\tauB(\phi)$ to quantify how geometry, receptor density, and intrinsic
kinetics shift the operational $C_{50}$.  Figure~4 evaluates the same
trajectory observables receptor by receptor on a sinusoidal surface, enabling
local curvature--retention and local $C_{50}$ maps.  Figure~5 uses the full
parametric response $\{C(\phi),\phi\}$ and the self/cross-rebinding
classification to compare dilute and nearly saturated regimes.

\section{Interpretive boundary}

The simulation establishes a kinetic mechanism: geometry changes the
probability of return to a receptor-bearing surface after dissociation, which
changes the cumulative bound time and therefore the concentration required to
maintain a specified occupancy.  The resulting $C_{50}$ is an operational
concentration scale for this open, diffusion-influenced kinetic model.  It
should not be described as a geometry-dependent thermodynamic $K_D$ unless an
independent equilibrium free-energy calculation demonstrates such a change.

\bibliographystyle{unsrtnat}
\bibliography{simulation_theory_references}

\end{document}
